\section{Introduction}

The isotropic constant compares the covariance of a uniform distribution
with the volume of its support.  For a convex body $K\subset\R^n$,
meaning a compact convex set with nonempty interior, write $|K|$ for its
Lebesgue volume, $b_K=|K|^{-1}\int_Kx\,\dd x$ for its centroid, and
\[
 \Sigma_K=\frac1{|K|}\int_K(x-b_K)(x-b_K)^{\mathsf T}\,\dd x.
\]
Its isotropic constant is
\begin{equation}\label{intro:isotropic-constant}
 L_K=\left(\frac{\det\Sigma_K}{|K|^2}\right)^{1/(2n)}.
\end{equation}
This quantity is unchanged by invertible affine maps.  The slicing
problem, arising in Bourgain's work on convex bodies
\cite{Bourgain1986}, asks whether it admits an upper bound independent
of dimension.  Klartag proved an $O(n^{1/4})$ bound
\cite{Klartag2006}.  Subsequent progress used stochastic localization,
introduced by Eldan \cite{Eldan2013}, leading through subpolynomial
and polylogarithmic estimates \cite{Chen2021,KlartagLehec2022,Klartag2023}
to Klartag and Lehec's dimension-independent bound using Guan's
covariance estimate \cite{Guan2024,KlartagLehec2025}.
The sharp form, also called the strong isotropic constant conjecture,
asks, in each dimension, whether a simplex maximizes $L_K$.
The distinction is substantial: a universal bound does not
identify the extremal value or an extremizing body.

In the plane, Campi, Colesanti, and Gronchi proved the simplex bound,
and Saroglou established uniqueness of the maximizing body
\cite{CampiColesantiGronchi1999,Saroglou2010}.
In higher dimensions, geometric variation arguments have imposed
restrictions on possible maximizers.  Rademacher proved that a
simplicial polytope maximizing the isotropic constant among all convex
bodies must be a simplex \cite{Rademacher2016}. More recent restrictions on local
maximizers in terms of decomposability appear in \cite{Kipp2026}.
We prove the sharp inequality for arbitrary convex bodies and classify
all equality cases.

\begin{theorem}[The simplex inequality]\label{intro:simplex}
For every integer $n\ge1$, every convex body $K\subset\R^n$ satisfies
the following bound, with equality if and only if $K$ is a simplex:
\begin{equation}\label{intro:simplex-bound}
 L_K\le
 \frac{(n!)^{1/n}}{(n+1)^{(n+1)/(2n)}\sqrt{n+2}}.
\end{equation}
\end{theorem}

The theorem also has a volume-product consequence.  For a convex body
$K$ with $0$ in its interior, define its polar by
$K^\circ=\{y\in\R^n:\langle x,y\rangle\le1\text{ for every }x\in K\}$.
Using exponential measures on convex cones, Klartag proved that the
sharp simplex bound implies
\begin{equation}\label{intro:mahler}
 |K|\,|K^\circ|\ge\frac{(n+1)^{n+1}}{(n!)^2},
\end{equation}
the nonsymmetric Mahler inequality
\cite[Corollary~1.2]{Klartag2018}.  Thus Theorem~\ref{intro:simplex}
also gives this inequality in every dimension.

The proof proceeds through differential entropy.  For a probability
density $f$ on $\R^m$, define, when the integrals are finite,
\[
 h(f)=-\int_{\R^m}f(x)\log f(x)\,\dd x,
 \qquad
 \Cov(f)=\int_{\R^m}(x-\bar x)(x-\bar x)^{\mathsf T}f(x)\,\dd x,
 \quad \bar x=\int_{\R^m}xf(x)\,\dd x.
\]
All logarithms are natural, and $0\log0=0$.  A log-concave density can
be represented, up to a null set, as $f=e^{-V}$ for a proper lower
semicontinuous convex function $V:\R^m\to(-\infty,+\infty]$.
The value $+\infty$ permits a density to vanish.  Such a probability
density has finite entropy and moments, and its covariance is positive
definite; the needed tail bound is proved in Lemma~\ref{geo:coercivity}.

\begin{theorem}[The sharp entropy bound]\label{intro:entropy}
Let $m\ge1$ and let $f$ be a log-concave probability density on $\R^m$.
Then
\begin{equation}\label{intro:entropy-bound}
 h(f)\ge m+\frac12\log\det\Cov(f).
\end{equation}
Equality holds if and only if $f$ is, almost everywhere, the density
of $M(E_1,\ldots,E_m)^{\mathsf T}+b$, where $M$ is an invertible real
$m\times m$ matrix, $b\in\R^m$, and the $E_i$ are independent random
variables with density $e^{-u}\mathbf1_{\{u\ge0\}}$.
\end{theorem}

Sharpness follows from a product of $m$ independent exponential variables
of mean one: its entropy is $m$ and its covariance is the identity.
Both sides of \eqref{intro:entropy-bound} increase by $\log|\det A|$
under an invertible affine change $x\mapsto Ax+b$.  Thus the theorem
asserts that a log-concave density with identity covariance has entropy
at least $m$.

Bobkov and Madiman developed the relation between entropy bounds for
log-concave measures and the slicing problem \cite{BobkovMadiman2011}.
Fradelizi and Mar\'in Sola formulated the sharp entropy statement
\eqref{intro:entropy-bound} and proved its equivalence, across dimensions,
with the simplex bound \cite[Conjecture~3(iii) and Proposition~5.2]{FradeliziMarinSola2026}.
In dimension one, Melbourne, Nayar, and Roberto proved the sharp bound,
attained by a one-sided exponential distribution
\cite[Theorem~1.1]{MelbourneNayarRoberto2026}.
Theorem~\ref{intro:entropy} establishes the inequality and identifies
all equality cases in every dimension.
The cone reduction used here follows the preceding geometric and
entropic work; we include it both to derive the dimensional constant
and to transfer the equality classification to convex bodies.

\subsection{Ideas of the proof}

We first assume that $f=e^{-V}$ is smooth and that the Hessian of $V$ is
bounded above and below by positive multiples of the identity.
Let $\nabla\phi$ transport standard Gaussian measure to $f$, and put
$J=D^2\phi$.  Brenier's theorem and Caffarelli's regularity and
contraction theory supply this positive definite matrix field and its
uniform bounds
\cite{Brenier1991,McCann1995,Caffarelli2000,CorderoErausquinFigalli2019}.
The scalar model is the transport from a Gaussian variable to a
mean-one exponential variable,
\[
 t(a)=-\log\Phi(-a),\qquad q(a)=t'(a),
\]
where $\Phi$ is the standard Gaussian distribution function.  Since $q$
is an increasing bijection onto $(0,\infty)$, spectral calculus defines
the symmetric matrix field $A=q^{-1}(J)$.

Two features of this parametrization drive the proof.  First, an affine
change of the target can be chosen so that $\E A=0$, where expectation
is Gaussian.  This nonlinear normalization is proved by a homotopy and
a compactness argument; ordinary covariance normalization does not
give the required identity.  Second, differentiating the transport
equation and using convexity of $V$ gives an energy estimate for $A$.
The estimate retains a nonnegative term involving pairs of distinct
eigenvalues of $A$.  Discarding that term would lose the control needed
for matrix functions whose derivatives do not commute.

The zero mean permits an orthogonal decomposition
\[
 A(x)=\sum_{r=1}^m x_rX_r+Y(x),
 \qquad X_r=\E[x_rA(x)],
\]
where $Y$ has Gaussian chaos degrees at least two.  The first term lies
in the equality space of the Gaussian Poincar\'e inequality and hence
has no energy surplus.  We instead use its constant matrices $X_r$ to
choose a supporting plane for the covariance log determinant.
The higher-degree part supplies a strict spectral gap.  A covariance
identity involving the resolvent of the Ornstein--Uhlenbeck operator
then expresses the signed deficit
$m+\tfrac12\log\det\Cov(f)-h(f)$ in terms of these two parts.

The resulting estimate has both noncommuting matrix terms and scalar
divided differences.  A commutator square controls the former.
One scalar inequality absorbs the cost of the supporting plane and
the terms without derivatives.  A second controls the divided
differences, charging their off-diagonal error to the
eigenvalue-separation term retained earlier.
After this cancellation, the entropy deficit is bounded above by a
quadratic expression that is nonpositive on every chaos of degree at
least two. Retaining strict slack also controls $Y$ and forces
$\sum_rX_r^2$ towards the identity when the entropy gap tends to zero.  All numerical constants in the scalar inequalities are
exact rationals.  Their verification uses polynomial bounds on a compact
interval and separate analytic estimates on both tails.

For equality, it is essential that this remainder applies to nearly
sharp densities. The expected extremizers have restricted support and
affine potentials, so equality cannot be studied solely within the
smooth, uniformly convex class. We approximate an arbitrary equality
density while preserving entropy and covariance. The remainder then
forces the reparametrized Jacobians to converge in Gaussian $L^2$ to
a linear field $A(x)=\sum_rx_rX_r$, with $\sum_rX_r^2=I$.
Since $q$ is Lipschitz in Frobenius norm, the original Jacobians
converge to $q(A)$, and Gaussian Poincar\'e gives convergence of the
centered maps.

The local compatibility of this limiting Jacobian is decisive.
The constant and linear terms of its mixed-derivative identities
force the symmetric matrices $X_r$ to commute. Simultaneous
orthogonal diagonalization then gives independent one-dimensional
exponential transports. Only the nonvanishing of $q'(0)$ and $q''(0)$
is needed for the commutation argument, isolating a potentially
reusable algebraic step. No uniform Hessian bounds on the smooth
approximants are required. Their affine normalizations are controlled
only after the limiting law and its covariance have been identified.

Section~\ref{tr:section} constructs the normalization and proves the
transport energy estimate.  Section~\ref{mat:section} decomposes the
entropy deficit and proves the smooth case using two scalar
propositions.  Section~\ref{sc:section} proves those propositions;
Appendix~\ref{cert:verification} supplies the quantitative one-variable
bounds.  Section~\ref{geo:section} removes the smoothness assumptions.
Section~\ref{sec:rigidity} classifies equality in the entropy bound.
Section~\ref{sec:geometry} applies that classification to the exponential
density on the cone over $K$, completing Theorem~\ref{intro:simplex}.
